%% ------------- English version ---------------
\documentclass[english,hidelinks]{sbrt}
\usepackage{newtxtext}
\usepackage[english]{babel}
\usepackage[utf8]{inputenc}
\newtheorem{theorem}{Theorem}
%% ---------------------------------------------
\setlength{\marginparwidth}{2cm}
\usepackage{todonotes}
\newcommand\rigel[1]{\todo[color=green, inline]{\textbf{Rigel}: #1}}
\usepackage{orcidlink}
\usepackage{hyperref}
\usepackage[
  compatibility=false,
  font=footnotesize,
  labelsep=period,
  format=plain,
  justification=raggedright,
  singlelinecheck=false
]{caption}

\begin{document}

\title{Embedded Systems-based Predictive Maintenance for Air-Conditioning Equipment}

\author{Igor Costa, Jorge Alves, Ian Dias and Davi Ito
  \thanks{The authors are with Ibmec, Rio de Janeiro, RJ, Brazil.
  E-mails: igorcosta.silva10a@gmail.com,
  jorgedocruzeiro07@gmail.com,
  iandpcarvalho@gmail.com,
  cdavi8249@gmail.com.}%
  }

\maketitle

\markboth{XLIV BRAZILIAN SYMPOSIUM ON TELECOMMUNICATIONS AND SIGNAL PROCESSING - SBrT 2026, SEPTEMBER 29TH TO OCTOBER 2ND, 2026, SALVADOR, BA}{}

\begin{abstract}
\end{abstract}
\begin{keywords}
Predictive maintenance, air-conditioning systems, embedded systems, Internet of Things, anomaly detection.
\end{keywords}

\section{Introduction}

Embedded systems integrate dedicated computation with sensors, actuators, and communication interfaces, enabling a physical asset to be monitored and controlled close to the point where data are produced. In the Internet of Things (IoT), this combination connects physical objects to networks and analytical services: Gubbi et al. describe the convergence of wireless sensor networks, the Internet, distributed computing, and embedded sensor and actuator nodes, while Lin et al. organize IoT architectures and discuss fog/edge processing near devices to reduce latency and improve service resilience. These characteristics make embedded IoT systems suitable for continuous equipment monitoring, because measurements can be acquired at the asset, processed locally or remotely, and made available for maintenance decisions \cite{gubbi2013iot,lin2017iot}.

Predictive maintenance uses condition and operational data to support maintenance decisions before an unexpected failure, combining data acquisition, processing, health assessment, and action planning. Carvalho et al. show in a systematic review that machine-learning applications in predictive maintenance depend on the characteristics and quality of the data, as well as on the selection and validation of the method. For HVAC systems, Es-Sakali et al. review knowledge-based, physics-based, and data-driven approaches that use IoT sensors and historical measurements to assess equipment health and anticipate failures. Motivated by this context, this project focuses on air-conditioning units that operate continuously, twenty-four hours a day, in environments where cooling cannot be interrupted without affecting essential activities, such as hospital wards, data centers and server rooms, and telecommunication shelters. Continuous duty is also convenient from a measurement standpoint, because the equipment remains in a steady operating regime, which allows a stable baseline to be established and prevents daily start-stop transients from being confused with progressive degradation. The proposed prototype will investigate continuous monitoring and early anomaly detection as a first step toward predictive maintenance, without claiming remaining-useful-life estimation before a representative history of validated failures is available \cite{carvalho2019mlpdm,essakali2022hvac}.

\textit{Related work.} Yousuf et al. developed an integrated condition-monitoring and fault-detection architecture for AC induction motors. Their system combines temperature, vibration, current, voltage, and speed sensors with an Arduino-based acquisition unit, local alarms, GSM notifications, relay-based protection, and historical storage on the Blynk IoT platform. Mohammed et al. proposed a complementary Industrial IoT architecture in which a Raspberry Pi collects vibration, current, and temperature measurements and transmits them to a cloud server through MQTT. These studies demonstrate how embedded acquisition and remote communication can provide continuous evidence about motor condition, although their results still require validation before being transferred to compressors and fan motors operating under the cycling and thermal-load variations of air-conditioning systems \cite{yousuf2024iot,mohammed2023iot}.

Kolok et al. investigated whether low-cost embedded processing could support predictive-maintenance experiments on small rotating machines. Their prototype uses an ESP32 with MEMS vibration and acoustic sensors, extracts time- and frequency-domain characteristics through RMS and FFT, and applies Isolation Forest to detect deviations from healthy operation. The system achieved an overall detection accuracy above 72.9\%, showing the feasibility of lightweight local analysis but also leaving room for improvement and equipment-specific calibration. Mohammed et al. compared five supervised algorithms on an AC-motor dataset containing normal operation and induced faults, with Random Forest obtaining the highest reported accuracy. Together, these works support the evaluation of lightweight anomaly-detection and classification methods, provided that the proposed project documents its sampling procedure, fault simulation, training data, and independent test data \cite{kolok2025lowcost,mohammed2023iot}.

At the system level, predictive maintenance depends not only on the selected model, but also on the organization and quality of the complete data pipeline. Meitz et al. organized 249 publications into a framework of nine categories and 73 attributes, covering condition monitoring, fault detection, degradation modeling, scheduling, data handling, prognostic techniques, and evaluation. Gupta et al. complemented this broad view with a field study of airport baggage-handling conveyors, where vibration data were collected from a live system despite the absence of historical run-to-failure records. Their pipeline distinguished anomalies from noisy observations, processed heterogeneous maintenance records, and compared four classifiers. Random Forest achieved the best reported motor--gearbox classification performance, with precision, recall, and F1-score of 0.86. These findings indicate that sensing, cleaning, labeling, and evaluation should be treated as connected parts of the proposed system rather than as isolated modeling tasks \cite{meitz2025pdmframework,gupta2023baggage}.

Beyond detecting that something deviates from normal operation, a maintenance system must indicate what deviated, so that the alert can be acted upon. Tormos et al. addressed this problem in air-conditioning units of urban bus fleets, where operational data are collected continuously but are rarely accompanied by reliable fault labels. The authors first map the available measurements onto the vapor-compression refrigeration cycle, which yields physically meaningful features and an operational baseline. An Isolation Forest then identifies deviations from that baseline, and Kernel SHAP attributes each deviation score to the variables that produced it, with the attributions grouped into physics-informed explanation families so that anomalies are read as refrigeration mechanisms rather than as isolated variables. The reported results distinguish physically plausible rare operating states, inspection candidates, and ambiguous cases without relying on confirmed fault labels. The study is based on bus fleet telematics rather than on stationary equipment, and the attribution step is computationally heavy for a microcontroller, which places it on the server side of the proposed architecture. It nevertheless supports two decisions of this project: learning the healthy baseline instead of requiring labelled failures, and reporting which measurement triggered each alert \cite{tormos2026hvacxai}.

The implementation, documentation, and experimental material produced in this project are publicly available in the project repository \cite{repo}.

This work proposes an embedded system that monitors the condition of air-conditioning units under uninterrupted duty and flags anomalies early. The prototype reads temperature, vibration, and electrical current at the equipment through a low-cost microcontroller and computes time- and frequency-domain features on the spot. The measurements then reach a remote application, which compares them against the healthy operating baseline and reports what deviated. Four aspects of this design are our contribution. The architecture pairs local acquisition and lightweight edge processing with remote communication, sized for equipment that stays in a steady regime and cannot have its cooling interrupted. The detection strategy learns the healthy baseline instead of requiring a labelled failure history, which is the realistic condition for equipment that has never been instrumented, and follows what Tormos et al. and Gupta et al. report from field deployments with the same absence of run-to-failure records \cite{tormos2026hvacxai,gupta2023baggage}. The system names the measurement that produced the alert, so the maintenance team gets something to act on instead of an undifferentiated warning. And the experimental protocol induces anomalies in a controlled way on a laboratory bench, keeping the data that establishes the baseline separate from the data used to evaluate it. We restrict the prototype to the condition monitoring, data handling, fault detection, and evaluation stages of the framework organized by Meitz et al. Degradation modelling, prognostics, and remaining-useful-life estimation stay out, because they need a representative history of validated failures that the prototype does not have \cite{meitz2025pdmframework}. The initial tests run on a bench built around a 12~V fan, the same class of machine as the condenser fan of an air-conditioning unit, which lets faults be induced in a controlled way. A bench fan is still not a compressor, so the results characterise the method rather than validate it on a real unit.

The rest of this work is organized as follows. Section~\ref{sec:system-description} describes the system under development, its sensors, microcontroller and communication protocol, together with the configuration defined for the first experiments. Section~\ref{sec:method} details the proposed method, covering signal acquisition, feature extraction, the operating baseline, the detection of deviations and the interpretation of each alert. We then report the results obtained with the prototype in Section~\ref{sec:results}, and Section~\ref{sec:conclusion} concludes the work, compares the proposed method with an alternative solution and outlines the next steps.

\rigel{Pendência da AC Semana 5: verificar as referências produzidas por alunos do Ibmec e citá-las caso seja oportuno.}

\section{System Description}
\label{sec:system-description}

The monitoring station is built around a single low-cost microcontroller that reads the sensors at the equipment, computes features locally, and sends the result onward. Kolok et al. showed that a board of this class can extract time- and frequency-domain characteristics and run a lightweight anomaly detector without a server, and Lin et al. describe why processing near the device reduces latency and keeps the service running when the link does not \cite{kolok2025lowcost,lin2017iot}. The prototype uses an ESP32 module with external PSRAM, which leaves room for the sample buffers that a frequency transform needs.

\subsection{Vibration Sensing}

A MEMS accelerometer is fixed to the shell of the rotating component and reports acceleration along three axes. The firmware computes the RMS of each window and a discrete Fourier transform, the same pair of characteristics that Kolok et al. applied to small rotating machines, and Yousuf et al. combine with electrical measurements on induction motors \cite{kolok2025lowcost,yousuf2024iot}. Mechanical coupling decides whether the reading means anything: a sensor that is not rigidly attached measures its own movement rather than the machine, so the accelerometer sits on a flat and rigid part of the casing, held by a magnetic base or structural adhesive.

\subsection{Electrical Current Sensing}

Current rises before most other evidence appears, because a motor under friction or obstruction draws more to hold its speed. Yousuf et al. and Mohammed et al. both include current among the quantities they acquire from motors under induced faults \cite{yousuf2024iot,mohammed2023iot}.

The prototype uses two different sensors for this single quantity. On the bench, where the machine runs at low voltage, a shunt monitor on the supply line measures the current directly. On an air-conditioning unit, where the compressor runs from the mains, a split-core current transformer clamps around one conductor and measures by induction, without cutting any wire and without contact with the supply. The distinction matters for safety as much as for accuracy, and it is declared here because the two setups are not interchangeable.

\subsection{Temperature Sensing}

A single temperature reading describes the room rather than the machine. The quantity that carries information is the difference between two points, which is why the station uses digital probes on a shared one-wire bus rather than one ambient sensor.

On the equipment, the probes are clamped to the liquid line and to the suction line of the refrigerant circuit. Tormos et al. map the available measurements onto the vapour-compression cycle precisely to obtain features with physical meaning, and Es-Sakali et al. review how HVAC health assessment rests on this kind of thermal evidence \cite{tormos2026hvacxai,essakali2022hvac}. Two further probes at the return and supply air of the indoor unit measure cooling capacity, and they have the practical advantage of requiring no access to the electrical enclosure.

\subsection{Processing and Communication}

The microcontroller acquires the three quantities, computes the window features, and publishes the result through MQTT, a lightweight publish and subscribe protocol for low-bandwidth and intermittent links. Mohammed et al. adopt the same arrangement, with an embedded acquisition unit transmitting vibration, current, and temperature to a remote server, and Gubbi et al. describe the wider architecture in which such devices deliver data to analytical services \cite{mohammed2023iot,gubbi2013iot}. The remote side compares each window against the healthy baseline and reports which measurement produced the alert.

\subsection{Experimental Configuration}

The first experiments run on a bench built around a 12~V fan driven by the same microcontroller, with the accelerometer on the motor body, the shunt monitor in series with the supply, and one probe on the motor and another measuring ambient air, so that heating is read as a rise above the room rather than as an absolute value.

Four conditions are induced in a controlled and repeatable way: a mass attached to one blade for imbalance, added friction on the shaft, a blocked air path for obstruction, and a loosened mounting screw. Each condition produces a different combination across the three quantities, which is what allows an alert to name the mechanism instead of only reporting that something changed.

Each condition is recorded in independent runs rather than sliced from a single acquisition. Gupta et al. report that ignoring this separation inflates the measured performance, and Meitz et al. treat evaluation as a stage of its own rather than an afterthought \cite{gupta2023baggage,meitz2025pdmframework}. The healthy baseline is built from runs that are never reused for evaluation.

\section{The Method}
\label{sec:method}

\rigel{AC Semana 7: Elaborar a seção do método proposto. Descreva o funcionamento do sistema, inclua um diagrama da arquitetura e, quando disponível, uma foto do protótipo. Um pseudocódigo também pode ser utilizado para apresentar os principais passos do método.}

\section{Results}
\label{sec:results}

\rigel{AC Semana 8: Elaborar a seção de resultados experimentais, apresentando os resultados qualitativos e quantitativos obtidos até o momento. Elaborar também o Abstract e a Conclusion.}

\rigel{AC Semana 9: Inserir os resultados quantitativos da avaliação, refinar o Abstract e concluir a discussão do desempenho do sistema embarcado proposto.}

\section{Conclusion}
\label{sec:conclusion}

\bibliographystyle{ieeetr}
\bibliography{refs}

\end{document}
